\subsection{Fragestellungen}
\label{subsec:fragestellungen}

Aus der beschriebenen Problemstellung und Zielsetzung ergeben sich sechs
zentrale Untersuchungsfragen. Sie strukturieren die Analyse des Templates und
bilden den Bezugsrahmen für seine spätere Bewertung:

\begin{enumerate}
  \item Welche fachlich-technischen, qualitativen und betrieblichen
    Anforderungen werden durch das Technologie-Template adressiert?
  \item Wie unterstützen die Architektur und das Profilmodell die
    Bereitstellung unterschiedlicher Varianten für Web-, Cloud- und
    Desktop-Projekte?
  \item Inwiefern fördern der gemeinsame technische Kern und die
    Single-Repository-Strategie die Wiederverwendung und Standardisierung
    zwischen Projekten?
  \item Welches Potenzial bietet das Template, wiederkehrende Aufwände bei
    Projektinitialisierung, Entwicklungsworkflow, Qualitätssicherung und
    Bereitstellung zu reduzieren?
  \item Für welche Projektarten und Kundenanforderungen stellt das Template
    eine geeignete technische Ausgangsbasis dar?
  \item Welche technischen Grenzen, Wartungsaufwände, Risiken und externen
    Abhängigkeiten sind bei seinem Einsatz zu berücksichtigen?
\end{enumerate}

Die Fragen sind miteinander verknüpft: Aussagen zu Produktivität und Eignung
setzen zunächst eine nachvollziehbare Betrachtung der Anforderungen, der
Architektur und der vorhandenen Verifikation voraus. Ihre Beantwortung wird
daher nicht isoliert, sondern entlang der aufeinander aufbauenden Kapitel der
Fallstudie vorgenommen.
